\clearpage
\begin{summarybox}
\textbf{一句话核心：} 先看镜面怎样改变横纵增量，再把像纵增量除以像横增量；实际反射保持长度，斜率计算只需要方向。

\medskip
\textbf{三条速查：}
\[
\boxed{\begin{array}{c|c}
\text{镜面}&\text{像斜率}\\ \hline
x=a\text{ 或 }y=b&-k\\
y=x+c\text{ 或 }y=-x+c&1/k\quad(k\ne0\text{ 且有限})
\end{array}}
\]
\[
\boxed{l:y=mx+c,\quad
k'=\frac{2m+(m^2-1)k}{1-m^2+2mk}}
\]
通式要求 $m,k$ 有限且分母非零；分母为零表示像直线竖直。原线或镜面竖直时改用增量判断。

\medskip
\textbf{使用顺序：}
\begin{enumerate}[leftmargin=*]
\item 辨镜面：水平、竖直、两种 $45^\circ$，还是任意斜线？
\item 判边界：像横增量是否为零？是则直接写竖直线。
\item 要方程再定位：找镜面上的公共点，或求原线上一点的对称点。
\item 遇内心或角平分线：继续核对射线朝向及角的内部。
\end{enumerate}

\medskip
\textbf{交卷前四查：}
\begin{itemize}[leftmargin=*]
\item 变号与倒数是否选对？是否误当成垂直条件？
\item 是否把“斜率不存在”误写成“像直线不存在”？
\item 方向吻合之外，位置是否也符合指定镜面？
\item 判断内角时，是否检查了射线，而非只看延长线？
\end{itemize}

\textbf{与 P018 的分工：} P055 先求像直线方向；需要具体点坐标时，再使用 P018 的点对称公式。
\end{summarybox}
